\documentclass[11pt, a4paper]{article}

\usepackage[utf8]{inputenc}
\usepackage[T1]{fontenc}
\usepackage[slovene]{babel}
\usepackage{lmodern}

\usepackage{xcolor}
\usepackage{bbold}
\usepackage{amsmath}
\usepackage{amssymb}

\input{../../../common/environments.tex}
\input{../../../common/settings.tex}

\input{../../../common/macroses/00_common_macroses.tex}
\input{../../../common/macroses/02_math_operators.tex}
\input{../../../common/macroses/03_shortcuts.tex}

\begin{document}

\title{Podatkovne strukture in algoritmi 2 -- Domača naloga 2}
\author{Ruslan Urazbakhtin}
\maketitle

\newpage

\section*{Hitra Fourierova transformacija}
\paragraph{1.\ naloga.} Dana sta niza \(t\) (dolžine \(n\)) in \(p\) (dolžine \(m\)). Poiščite vse indekse \(i\), za katere velja \(t[i : i + m] = p\).
%
\paragraph{Rešitev.} Naj bo \(\Gamma\) dana abeceda. Za začetek recimo, da je \(\Gamma = \set{a, b}\). Najprej transformiramo abecedo na naslednji način
\[
    a \mapsto 1, \qquad b \mapsto -1.
\]
Priredimo nizu \(t\) polinom \(t(X)\) na naslednji način:
\[
    t(X)[j] = t[j].
\]
Priredimo nizu \(p\) polinom \(p(X)\) na naslednji način:
\[
    p(X)[j] = p[m - 1 - j],
\]
torej obrnemo niz. Oglejmo si splošni koeficient pri \(X^k\) v produktu \(q(X) = t(X) \cdot p(X)\):
\[
    q[k] = t_0p_{m-1-k} + t_1p_{m-k} + \cdots + t_kp_{m-1} = \sum_{j=0}^{k} t_j p_{m-1-k + j},
\]
kjer je \(p_j = 0\), če \(j < 0\) ali \(j \geq m\).
%
Zdaj pa pogledamo konkretne vrednosti \(k\):
%
\begin{align*}
    &k = m - 1 \lthen q[k] = \sum_{j=0}^{m-1} t_jp_j = t[0 : m] \cdot p \\
    &k = m \lthen q[k] = \sum_{j=0}^{m} t_j p_{j - 1} = \sum_{j=0}^{m-1} t_{j+1}p_j =  t[1 : m + 1] \cdot p\\
    &\vdots\\
    &k = m - 1 + i \lthen q[k] = \sum_{j=0}^{m-1+i} t_j p_{j - i} = \sum_{j=0}^{m-1} t_{j+i}p_j =  t[i : m + i] \cdot p,
\end{align*}
kjer \(t[i : m + i] \cdot p\) označuje skalarni produkt vektorjev \(t[i : m + i]\) in \(p\). Vidimo, da za \(i \in \set{0, \ldots, n - m}\) pri potenci \(X^{m - 1 + i}\) dobivamo ravno skalarne produkte vektorjev \(t[i : m + i]\) in \(p\). Preden zapišemo algoritem, dokažemo trditev.

\begin{trditev}
    \label{proof-1}
    Naj bosta \(\vec{a}, \vec{b} \in \set{-1, 1}^n\). Tedaj 
    \[
        \vec{a} = \vec{b} \liff \vec{a} \cdot \vec{b} = n.
    \]
\end{trditev}

\begin{proof}
    \((\Rightarrow)\) Recimo, da je \(\vec{a} = \vec{b}\). Sledi, da 
    \[
        \vec{a} \cdot \vec{b} = \vec{a} \cdot \vec{a} = \norm{\vec{a}}^2 = n.
    \]
    \((\Leftarrow)\) Recimo, da je \(\vec{a} \cdot \vec{b} = n\). Sledi, da 
    \[
        n = \sum_{j=1}^{n}a_jb_j = k - l,
    \]
    kjer je \(k = |\setb{j \in [n]}{a_jb_j = 1}|\) in \(l = |\setb{j \in [n]}{a_jb_j = -1}|\). Ker je \(k \leq n\), enakost je dosežena le v primeru, ko \(k = n\) in \(l = 0\) oziroma, ko 
    \[
        \all{j \in [n]} a_j b_j = 1 \liff \all{j \in [n]} a_j = b_j. \qedhere
    \]
\end{proof}
\noindent
Torej, dovolj je poiskati vse indekse \(i \in \set{0, \ldots, n - m}\), da \(t(X)[i : m + i] \cdot p(X) = m\).

\paragraph{Algoritem.}
\begin{enumerate}
    \item Priredi nizoma \(t\) in \(p\) polinoma \(t(X)\) in \(p(X)\). Časovna zahtevnost je \(O(n)\).
    \item Izračunaj \(q(X) = \text{IFFT}(\text{FFT}(t(X)) \cdot \text{FFT}(p(X)))\). Časovna zahtevnost je \(O(n \log n)\).
    \item Poišči \(i \in \set{0, \ldots, n - m}\), da \(q[m - 1 + i] = m\). Časovna zahtevnost je \(O(n)\).
\end{enumerate}
\
\\
Za splošno abecedo \(\Gamma = \set{a_0, \ldots, a_{N - 1}}\) pa naredimo transformacijo 
\[
    \all{k \in \set{0, 1, \ldots, N-1}} a_k \mapsto e^{\frac{2 \pi i}{N} k}.
\]
Priredimo nizu \(t\) polinom \(t(X)\) na naslednji način:
\[
    t(X)[j] = t[j].
\]
Priredimo nizu \(p\) polinom \(p(X)\) na naslednji način:
\[
    p(X)[j] = \overline{p[m - 1 - j]},
\]
torej obrnemo niz in transponiramo njegove vrednosti. Spet za \(i \in \set{0, \ldots, n - m}\) pri potenci \(X^{m - 1 + i}\) dobivamo ravno skalarne produkte vektorjev \(t[i : m + i]\) in \(p\). Uporabimo isti algoritem kot prej, ki je pravilen tudi v splošnem primeru, saj velja

\begin{lema}
    Naj bosta \(\vec{a}, \vec{b} \in \C^n\) in \(\all{j \in [n]} |a_j| = |b_j| = 1\). Tedaj
    \[
        \vec{a} = \vec{b} \liff  \all{j \in [n]} a_j \overline{b_j} = 1.
    \] 
\end{lema}

\begin{proof}
    \((\Rightarrow)\) Recimo, da je \(\vec{a} = \vec{b}\). Naj bo \(j \in [n]\). Velja
    \[
        a_j\overline{b_j} = a_j \overline{a_j} = |a_j|^2 = 1.
    \]
    \((\Leftarrow)\) Recimo, da je \(\all{j \in [n]} a_j \overline{b_j} = 1\). Naj bo \(j \in [n]\). Označimo \(a_j = e^{i\phi_j}\) in \(b_j = e^{i \psi_j}\). Sledi, da
    \[
        1 = a_j \overline{b_j} = e^{i\phi_j}e^{-i \psi_j} = e^{i (\phi_j - \psi_j)} \lthen \phi_j - \psi_j = 2 \pi k \lthen \phi_j = \psi_j + 2 \pi k.
    \]
    Torej 
    \[
        a_j = e^{i\phi_j} = e^{i(\psi_j + 2 \pi k)} = e^{i\psi_j}e^{i2 \pi k} = e^{i\psi_j} = b_j. \qedhere
    \]
\end{proof}

\begin{trditev}
    Naj bosta \(\vec{a}, \vec{b} \in \C^n\) in \(\all{j \in [n]} |a_j| = |b_j| = 1\). Tedaj 
    \[
        \vec{a} = \vec{b} \liff \vec{a} \cdot \vec{b} = n.
    \]
\end{trditev}

\begin{proof}
    \((\Rightarrow)\) Recimo, da je \(\vec{a} = \vec{b}\). Tedaj 
    \[
        \vec{a} \cdot \vec{b} = \vec{a} \cdot \vec{a} = \sum_{j=1}^{n} a_j \overline{a_j} = \sum_{j=1}^{n} |a_j|^2 = \sum_{j=1}^{n} 1 = n.
    \]
    \((\Leftarrow)\) Recimo, da je \(\vec{a} \cdot \vec{b} = n\). Označimo 
    \[
        \all{j \in [n]} a_j = \cos \phi_j + i \sin \phi_j \quad \text{in} \quad \all{j \in [n]} b_j = \cos \psi_j + i \sin \psi_j.
    \]
    Tedaj 
    \begin{multline*}
        n = \vec{a} \cdot \vec{b} = \sum_{j=1}^{n} a_j \overline{b_j} = \sum_{j=1}^{n} (\cos \phi_j + i \sin \phi_j)(\cos \psi_j - i \sin \psi_j) = \\ = \sum_{j=1}^{n} (\cos (\phi_j - \psi_j) + i \sin (\phi_j - \psi_j)).
    \end{multline*}
    Torej 
    \[
        n = \sum_{j=1}^{n} \cos (\phi_j - \psi_j) \lthen \all{j \in [n]} \cos(\phi_j - \psi_j) = 1 \lthen \all{j \in [n]} \phi_j = \psi_j + 2 \pi k.
    \]
    Torej 
    \[
        \all{j \in [n]} a_j\overline{b_j} = 1 \lthen \vec{a} = \vec{b}. \qedhere
    \]
\end{proof}
%
\paragraph{2.\ naloga.} Zapišite algoritem za množenje dveh dolgih števil, ki ju predstavimo kot seznama števk (dolžine \(m\) in dolžine \(n\)), ki naj ima kompleksnost \(O((m+n) \log(m+n))\).
%
\paragraph{Rešitev.} Vsako število lahko predstavimo v obliki 
\[
    a = a_0 + a_1 X + \ldots + a_{m - 1} X^{m - 1}, \qquad b = b_0 + b_1 X + \ldots + b_{n - 1} X^{n - 1},
\]
kjer je \(X\) neka izbrana baza, npr.\ \(10\). Priredimo vsakemu številu polinom 
\[
    \all{j \in [n + m]} p(X)[j] = p[j],
\]
pri čemer manjkajoče indekse štejemo za 0. Tako dobimo polinoma \(a(X)\) in \(b(X)\). Tedaj množenje števil \(a\) in \(b\) ustreza množenju ustreznih polinomov ter prenosu prepolnjenih razredov.

\paragraph{Algoritem.}
\begin{enumerate}
    \item Priredi številoma \(a\) in \(b\) polinoma \(a(X)\) in \(b(X)\). Časovna zahtevnost je \(O(n + m)\).
    \item Izračunaj \(c(X) = \text{IFFT}(\text{FFT}(a(X)) \cdot \text{FFT}(b(X)))\).\\ Časovna zahtevnost je \(O((m+n) \log (m+n))\).
    \item Normaliziruj \(c(X)\), tj.
    \begin{enumerate}
        \item \(\text{carry := 0}\)
        \item Za vsak \(i \in \set{0, 1, \ldots, n + m - 1}\) ponavljaj
        \begin{enumerate}
            \item \(c(X)[i] = c(X)[i] + \text{carry}\);
            \item \(\text{carry} = c(X)[i]\, /\, 10\);
            \item \(c(X)[i] = c(X)[i] \mod 10\).
        \end{enumerate}
        \item Če \(\text{carry} \neq 0\), \(c(X)[m+n] = \text{carry}\).
        \item Pobriši vodilne ničle.
    \end{enumerate}
    Časovna zahtevnost je \(O(n + m)\).
\end{enumerate}

\end{document}
